\chapter*{Abstract}
\chaptermark{Abstract}
\addcontentsline{toc}{chapter}{Abstract}

\noindent\textbf{Titel:} Een Schaalbare Multi-Agent\-architectuur voor Enterprise Search\par
\noindent\textbf{Student:} Arne Albrecht\par
\noindent\textbf{Opleiding:} Master in de industriële wetenschappen: informatica\par
\noindent\textbf{Bedrijfspartner:} Easi\par
\noindent\textbf{Promotor(en) en begeleider(s):} Prof. dr. ir. Filip De Turck, dhr. Joris Steppe (Easi), dhr. Dorian Féaux (Easi)

\vspace{1em}

\noindent Binnen moderne organisaties bevindt relevante informatie zich verspreid over verschillende systemen, zoals e-mailomgevingen, documentplatformen en CRM-toepassingen. Daardoor wordt het steeds moeilijker om snel een volledig en bruikbaar antwoord op een informatievraag te formuleren. Deze masterproef onderzoekt hoe een multi-agent\-architectuur kan worden ingezet om bedrijfsdata uit meerdere systemen op een veilige en efficiënte manier te doorzoeken. Hiervoor werd een prototype ontwikkeld met een centrale orchestrator en drie gespecialiseerde subagents voor Microsoft Graph, Salesforce en SmartSales. De communicatie met de onderliggende databronnen verloopt via het Model Context Protocol (MCP), terwijl Microsoft Agent Framework en Azure OpenAI gebruikt werden voor de agentlogica. Het ontwikkelde prototype toont aan dat deze architectuur technisch haalbaar is en dat informatie uit meerdere systemen op een gecontroleerde en uitbreidbare manier kan worden samengebracht. De evaluatie gebeurde aan de hand van testqueries waarbij onder meer correctheid van antwoorden, query routing, responstijd en tokenverbruik werden onderzocht. De eerste resultaten geven aan dat de voorgestelde aanpak potentieel heeft voor enterprise search binnen bedrijfsomgevingen, maar tonen ook aan dat de uiteindelijke kwaliteit sterk afhangt van de betrouwbaarheid en maturiteit van de onderliggende agent- en toolimplementaties. Deze masterproef levert daarmee een proof-of-concept en een basis voor verdere verfijning en uitbreiding in een volgende ontwikkelfase.

\vspace{1em}

\noindent\textbf{Trefwoorden:} multi-agent\-architectuur, enterprise search, Agentic AI, Model Context Protocol, Large Language Model